% theory_dynamic_traffic_power.tex -- EV/traffic/power general theory
\section{动态交通流、均衡与电网耦合理论}\label{sec:evpt-theory}

\begin{theorynote}
本章给出 CTM、LTM、动态用户均衡、社会福利优化与充电电网耦合的通用理论。
它用于推导守恒量、离散稳定条件和最优性证书；A--H 各入口实际使用的近似、变量与
拒绝分支以源码等价章为准，不能把本章的统一表述解释为一个已投产的全局模型。
\end{theorynote}

\subsection{三角基本图与 CTM 守恒}
对路段 $a$，密度 $k_a$ 的三角基本图为
\begin{equation}
 q_a(k)=\min\{v_a^f k,\bar q_a,w_a(k_a^{jam}-k)\}.
\end{equation}
把路段离散为长度 $\delta_{am}$ 的元胞，时间步长为 $\Delta t$。发送量和接收量为
\begin{align}
 S_{am}^n&=\Delta t\min\{v_a^f n_{am}^n/\delta_{am},\bar q_a\},\\
 R_{am}^n&=\Delta t\min\{\bar q_a,w_a(N_{am}^{jam}-n_{am}^n)/\delta_{am}\},
\end{align}
边界通量 $y_{am}^n=\min(S_{am}^n,R_{a,m+1}^n)$，于是
\begin{equation}
 n_{am}^{n+1}=n_{am}^n+y_{a,m-1}^n-y_{am}^n.
\end{equation}
对所有元胞求和后内部边界通量相消，因此总车辆变化只等于边界流入减流出。这是数值
准入中必须独立重算的守恒式。

单调 CTM 的 CFL 条件至少要求
\begin{equation}
 \Delta t\le\min_{a,m}\left\{\frac{\delta_{am}}{v_a^f},
 \frac{\delta_{am}}{w_a}\right\}.
\end{equation}
违反 CFL 时，即使程序返回有限数值，也不能解释为物理可信的传播结果。

\subsection{LTM 累积曲线}
LTM 用累积进入/离开曲线 $N_a^{in}(n)$、$N_a^{out}(n)$ 代替空间元胞。离散自由流与
反向波延迟为
\begin{equation}
 \tau_a^{ff}=\left\lceil\frac{L_a}{v_a^f\Delta t}\right\rceil,
 \qquad
 \tau_a^{bw}=\left\lceil\frac{L_a}{w_a\Delta t}\right\rceil.
\end{equation}
发送和接收上界为
\begin{align}
 S_a^n&=\min\{N_a^{in}(n-\tau_a^{ff})-N_a^{out}(n),\bar q_a\Delta t\}_+,\\
 R_a^n&=\min\{N_a^{out}(n-\tau_a^{bw})+N_a^{jam}-N_a^{in}(n),
                    \bar q_a\Delta t\}_+.
\end{align}
CTM 与 LTM 共享连续守恒理论，但有限步长下存在不同启动延迟和容量偏差；二者相互比较是
离散模型交叉诊断，不是“任一结果自动成为真值”。

\subsection{Wardrop 均衡与变分不等式}
令 $\mathcal X$ 为满足 OD 守恒和非负性的路径流集合，$c(x)$ 为交通传播、排队、充电价格
共同生成的广义成本。动态用户均衡满足变分不等式
\begin{equation}
 c(x^*)^\mathsf T(x-x^*)\ge0,\qquad \forall x\in\mathcal X.
\end{equation}
若成本可积且 Jacobian 对称半正定，该问题等价于 Beckmann 势函数最小化；一般动态交通
和价格反馈不保证这些条件。MSA 更新
\begin{equation}
 x^{k+1}=x^k+\alpha_k(y^k-x^k),\qquad \alpha_k=\frac1{k+1},
\end{equation}
只在适当单调性和精确全有或全无子问题下具有经典收敛保证。报告的 relative gap/VI gap
必须与全局最优性证书分开。

\subsection{充电站库存与电网功率}
站点等待、服务和离开人口满足离散守恒，例如
\begin{align}
 Q_{s,n+1}&=Q_{s,n}+u_{s,n}-b_{s,n},\\
 B_{s,n+1}&=B_{s,n}+b_{s,n}-c_{s,n},\\
 H_{s,n+1}&=H_{s,n}+c_{s,n}-o_{s,n}.
\end{align}
若车辆 $v$ 以 $p_{v,n}^{ch}$ 充电、以 $p_{v,n}^{dis}$ 放电，则电量递推为
\begin{equation}
 E_{v,n+1}=E_{v,n}+\eta_v^{ch}p_{v,n}^{ch}\Delta t
 -p_{v,n}^{dis}\Delta t/\eta_v^{dis}-e_{v,n}^{drive}.
\end{equation}
站点对电网的有功注入按 kW 到 MW 换算；价格从 \texttt{\$/MWh} 换算到
\texttt{\$/kWh} 时除以 1000。
交通车辆数、站点 kW 与电网 MW 三种量纲不得直接相加。

\subsection{社会福利、MPEC 与证书层级}
联合规划的通式可写为
\begin{equation}
 \max\;B_{travel}-C_{time}-C_{queue}-C_{generation}-C_{infrastructure},
\end{equation}
并同时满足交通守恒、充电库存和电网平衡。Wardrop 条件可作为 VI、互补约束或有限路径
PWL/MILP 近似嵌入。LP 的最优状态、MILP 的全局 gap、NLP 的局部 KKT 残差和启发式固定点
属于不同证据层级，不能用同一个 \fld{converged} 混合表达。

\begin{gapnote}
当前没有一个入口同时对连续动态交通、非线性 AC/DC 电网、离散站点建设、个体 SOC 与
全局 Wardrop 互补给出统一全局证书。认证型 MILP 只认证其有限路径和 PWL 模型；NLP 只提供
局部驻点；顶层 A 为分解仿真。该边界由结果中的模型验证和证书字段保持可见。
\end{gapnote}

\subsection{与实现的对应}
\begin{center}\footnotesize
\begin{longtable}{@{}L{6.2cm}L{9.0cm}@{}}
\toprule 理论条目 & 实现状态\\\midrule
\endfirsthead\toprule 理论条目 & 实现状态\\\midrule\endhead
CTM 三角基本图与守恒 & \implfull{src/ev\_power\_traffic/ctm\_propagation.cpp:ctm\_forward\_pass}\\
LTM 累积曲线与站点库存 & \implfull{src/ev\_power\_traffic/ltm\_propagation.cpp:simulate\_ltm\_network}\\
DUE--VI 与 MSA/线搜索 & \implfull{src/ev\_power\_traffic/ctm\_due\_vi.cpp:solve\_ctm\_due\_vi}\\
价格协调与 DCOPF & \implpart{分解迭代；只认证价格不动点，不是统一社会福利最优}\\
有限模型 LP/MILP/MPEC & \implfull{src/ev\_power\_traffic/joint\_optimizer.cpp:solve\_joint\_optimizer}\\
全连续 AC/DC--交通全局证书 & \implnone\\
\bottomrule
\end{longtable}
\end{center}
